% ATTEMPT_STATUS: unresolved
% TOP_PROBLEM: {"problemId": "problem.steins-conjecture-on-the-hilbert-transform-along-lipschitz-vector-fields", "canonicalTitle": "Stein's Conjecture on the Hilbert Transform along Lipschitz Vector Fields", "releaseRank": 434, "primaryDomain": "analysis_pde", "primaryDomainLabel": "Analysis & PDE", "displayStatus": "open_with_solved_subcases", "releaseStatus": "open_with_solved_subcases"}
% !TeX program = lualatex
\documentclass[11pt]{article}
\usepackage[margin=25mm]{geometry}
\usepackage{fontspec}
\setmainfont{Latin Modern Roman}
\usepackage{amsmath,amssymb,mathtools}
\usepackage{unicode-math}
\setmathfont{Latin Modern Math}
\usepackage{xurl,hyperref}
\hypersetup{colorlinks=true,urlcolor=blue}
\setcounter{secnumdepth}{0}
\setlength{\parindent}{0pt}
\setlength{\parskip}{5pt}
\setlength{\emergencystretch}{3em}
\begin{document}
\section*{\texorpdfstring{Stein's Conjecture on the Hilbert Transform along Lipschitz Vector Fields}{Stein's Conjecture on the Hilbert Transform along Lipschitz Vector Fields}}\label{434}
\subsection{Definitions and mathematical statement}
Let $v:{\mathbb{R}}^2\to S^1$ be Lipschitz, with constant $L$, and define the locally truncated directional Hilbert transform on smooth compactly supported functions by
\[
 H_{v,\epsilon}f(x)=\lim_{\delta\downarrow0}
 \int_{\delta<|t|<\epsilon}f(x-tv(x))\,\frac{{\,\mathrm{d}} t}{t}.
\]
The direction is evaluated at the base point $x$, not followed along an integral curve. The question asks for $L^2$ boundedness at a truncation scale controlled by the Lipschitz scale, typically $\epsilon$ a sufficiently small multiple of $L^{-1}$ when $L>0$:
\[
 \|H_{v,\epsilon}f\|_2\le C_{v,\epsilon}\|f\|_2.
\]
The operator is understood initially through symmetric truncations and then its bounded extension. The catalogue does not explicitly fix a uniform numerical operator-norm constant; no such numerical value is imposed here. Constant vector fields are treated separately from the expression $L^{-1}$.
\par\smallskip\textit{Formulation and concept references:} \href{https://doi.org/10.1090/S0065-9266-10-00572-7}{[S1]}.

\subsection{Short English statement}
Does a Hilbert transform remain bounded on square-integrable functions when its direction varies in a Lipschitz manner and the integration is suitably local?
\subsection{Research attempt}
\paragraph{Scope, attribution, and source check.}
Prepared by GPT-6 Astra Ultra on 27 September 2026, this unresolved
attempt uses the base-point direction in the statement and symmetric
truncation at zero. The integral is not along a flow trajectory.
The Lacey--Li author version [E1] of [S1] explicitly formulates the
Lipschitz conjecture. Its maximal-function and additional-regularity
results are not substituted for the requested Hilbert-transform
bound. The work here proves several elementary estimates and
isolates the cancellation estimate they fail to supply.
For definiteness take $0<\epsilon L<1/2$ when $L>0$.

\paragraph{The constant-direction case.}
For $v(x)=v_0$, Fourier transformation diagonalizes the truncated
operator. With Fourier phases $e^{-ix\cdot\xi}$, its multiplier is
\[
 m_{\delta,\epsilon}(\xi)
 =-2i\int_\delta^\epsilon
                    \frac{\sin(t\,\xi\cdot v_0)}t\,dt.
\]
The primitive $\operatorname{Si}(R)=\int_0^R\sin u/u\,du$ is
uniformly bounded. For example the integral on $[0,1]$ is at
most one in absolute value; integration by parts bounds every
tail beginning at one by a fixed constant, independent of its
upper endpoint. The multiplier therefore has a uniform bound
independent of $\delta,\epsilon,\xi$. Plancherel proves an
$L^2$ bound, and dominated convergence of multipliers gives
a bounded strong limit as $\delta\downarrow0$.

The cancellation of the odd kernel is essential. The absolute
kernel has integral $2\log(\epsilon/\delta)$, which diverges.
A proof by absolute values alone cannot even recover this
constant-field example.

\paragraph{The geometry of each short displacement.}
Define $\Phi_t(x)=x-tv(x)$. Lipschitz continuity gives
\[
 (1-|t|L)|x-y|\leq|\Phi_t(x)-\Phi_t(y)|
                         \leq(1+|t|L)|x-y|.
\]
When $|t|L<1$, it is a global bijection: solving
$x=z+tv(x)$ is a contraction problem for every $z$.
Rademacher differentiability and the change-of-variables
formula for bi-Lipschitz maps imply
$|\det D\Phi_t|\geq(1-|t|L)^2$ almost everywhere. Hence
\begin{equation}\label{a434:composition}
 \|f\circ\Phi_t\|_2\leq(1-|t|L)^{-1}\|f\|_2.
\end{equation}
This estimate controls the geometry and rules out a fold
at the chosen scale. It does not control cancellation
between different values of $t$.

Pairing positive and negative $t$ and applying the fundamental
theorem of calculus along the straight segment gives
\[
 f(x-tv(x))-f(x+tv(x))
 =-\int_{-t}^t\nabla f(x-sv(x))\cdot v(x)\,ds.
\]
Combining this with \eqref{a434:composition} proves
\begin{equation}\label{a434:h1}
 \|H_{v,\epsilon}f\|_2
 \leq\frac{2\epsilon}{1-\epsilon L}\|\nabla f\|_2
 \qquad(f\in H^1(\mathbb R^2)).
\end{equation}
It also proves convergence of the small-scale tail on $H^1$,
since replacing $\epsilon$ by a smaller upper endpoint
makes the right side tend to zero. Thus there is an actual
dense-domain definition, rather than just a formal
principal-value notation.

\paragraph{A scaling check.}
The dimensionless small parameter is $\epsilon L$.
In two dimensions $D_af(x)=a f(ax)$ is unitary on $L^2$.
A change of variable in the integral gives
\[
 D_a^{-1}H_{v,\epsilon}D_a
       =H_{v_a,a\epsilon},\qquad v_a(y)=v(y/a).
\]
The new field has Lipschitz constant $L/a$, so the product
of that constant and the new truncation is still
$\epsilon L$. Taking $a=L$ normalizes a positive
Lipschitz constant to one. It does not improve the
small parameter or eliminate high frequencies.
Constant fields, where $L=0$, are handled by the
separate Fourier calculation rather than this dilation.

\paragraph{Why this is not an $L^2$ proof.}
For a function whose Fourier support lies in $|\xi|\leq R$,
\eqref{a434:h1} gives a bound of size $O(\epsilon R)$.
The desired operator norm cannot grow with the input's
frequency. Taking $f(x)=e^{iR\xi_0\cdot x}\chi(x)$ makes
$\|\nabla f\|_2$ grow like $R\|f\|_2$ while the conjectured
bound remains independent of $R$. Smooth compact support
does not remove arbitrarily high frequencies.

The same issue obstructs freezing the field on small patches.
If $v$ is uniformly close to $v_0$, interpolating their
straight displacements bounds the difference of their
actions on $H^1$ by a constant times
$\epsilon\|v-v_0\|_\infty\|\nabla f\|_2$, provided the
interpolated displacement maps remain bi-Lipschitz.
Closeness of directions is multiplied by the frequency.
No fixed spatial partition makes that error small on
all of $L^2$. A frequency-dependent decomposition would
need further orthogonality estimates between its pieces.

\paragraph{The exact summation problem over scales.}
Let $\epsilon_j=2^{-j}\epsilon$ and define
\[
 T_jf(x)=\int_{\epsilon_{j+1}<|t|<\epsilon_j}
                            f(x-tv(x))\,\frac{dt}{t}.
\]
By \eqref{a434:composition},
\[
 \|T_j\|_{2\to2}\leq
                 \frac{2\log2}{1-\epsilon L}
\]
uniformly in $j$. The full transform is the strong
$H^1\to L^2$ sum of these operators. Summing their uniform
$L^2$ norms, however, produces an infinite series.

A precise sufficient missing estimate is two-sided
almost orthogonality: for constants $A,\alpha>0$,
\begin{equation}\label{a434:cotlar}
 \|T_j^*T_k\|_{2\to2}+\|T_jT_k^*\|_{2\to2}
                         \leq A^2\,2^{-\alpha|j-k|}.
\end{equation}
The Cotlar--Stein summation lemma, in its usual form with
summable square roots of these bounds, then gives a
uniform $L^2$ bound for every finite partial sum. The
sum of $2^{-\alpha|j-k|/2}$ over the difference of
indices is finite. Uniform boundedness plus convergence
on the dense subspace $H^1$ extends the limit to all
of $L^2$. This last step follows by approximating an
arbitrary input by an $H^1$ input and using the common
operator norm on the difference.

Condition \eqref{a434:cotlar} is explicitly an additional
hypothesis, not a theorem deduced here from Lipschitz
regularity. In particular the adjoint includes the
inverse displacement map and its Jacobian; treating
$T_j^*$ as simply $-T_j$ is generally incorrect for a
variable field. The operators at distinct base points
sample different directions, which is precisely what
the constant-direction Fourier calculation avoids.

\paragraph{Straightening a flow changes the operator.}
Let $\Psi_{-t}(x)$ denote the trajectory at time $-t$
of the vector field. Since $|v|=1$,
$|\Psi_{-s}(x)-x|\leq s$. The integral equation gives
\[
 |\Psi_{-t}(x)-(x-tv(x))|
 \leq\int_0^t L|\Psi_{-s}(x)-x|\,ds
 \leq Lt^2/2.
\]
This second-order positional error looks harmless
until it is applied to high-frequency functions:
an oscillation of frequency $R$ can acquire phase
error of order $RLt^2$. There is no bound uniform
in $R$ from this estimate alone. Therefore a
coordinate chart that straightens integral curves
does not conjugate the operator in the catalogue
to a constant-direction Hilbert transform without
a nontrivial error analysis.

This test is compatible with the proved $H^1$
estimate; both lose derivatives. It does not
construct a Lipschitz field for which the target
operator is unbounded. It rules out an unjustified
identification of two different singular integrals.

\paragraph{Verification and remaining gap.}
The Jacobian power in \eqref{a434:composition} is
dimension two, and taking its square root gives
the stated $L^2$ constant. The scale condition
keeps all intermediate maps invertible. The
dyadic integrals use both signs of $t$, and the
constant-field multiplier has the same convention.
The small-scale tail estimate proves convergence
before any conditional summation lemma is invoked.

The first missing step is a cancellation or
frequency-space estimate replacing the divergent
sum of single-scale norms, for the full Lipschitz
class at a suitably small scale. Proving
\eqref{a434:cotlar} would be one sufficient route;
it is not claimed to be necessary or to be the
best formulation of a successful proof.
\textbf{UNRESOLVED.} The constant-field theorem,
bi-Lipschitz geometry and dense-domain bounds are
proved, but no $L^2$ estimate for arbitrary
Lipschitz fields is obtained.

\subsection{Sources}
\begin{itemize}
\item[S1]On a Conjecture of E. M. Stein on the Hilbert Transform on Vector Fields.
\url{https://doi.org/10.1090/S0065-9266-10-00572-7}.
\textit{Formulation source.}
\item[E1]Michael Lacey and Xiaochun Li, \emph{On a Conjecture of E. M. Stein on the Hilbert Transform on Vector Fields}, author version of [S1].
\url{https://arxiv.org/abs/0704.0808}.
\textit{Primary source checked 27 September 2026 for the Lipschitz formulation and the scope of the partial results.}
\end{itemize}
\end{document}
